% Compile with XeLaTeX or LuaLaTeX because this report uses Mongolian Cyrillic.
\documentclass[12pt,a4paper]{article}

\usepackage{fontspec}
\IfFontExistsTF{Times New Roman}
  {\setmainfont{Times New Roman}}
  {\setmainfont{DejaVu Serif}}
\IfFontExistsTF{Arial}
  {\setsansfont{Arial}}
  {\setsansfont{DejaVu Sans}}

\usepackage[a4paper,margin=2.5cm]{geometry}
\usepackage{setspace}
\usepackage{enumitem}
\usepackage{array}
\usepackage{booktabs}
\usepackage{longtable}
\usepackage{xcolor}
\usepackage{hyperref}
\usepackage{listings}
\usepackage{caption}

\hypersetup{
  unicode=true,
  colorlinks=true,
  linkcolor=black,
  urlcolor=blue
}

\setstretch{1.15}
\setlist[itemize]{leftmargin=1.4cm}
\setlist[enumerate]{leftmargin=1.4cm}

\lstdefinestyle{api}{
  basicstyle=\ttfamily\small,
  breaklines=true,
  frame=single,
  columns=fullflexible
}

\newcommand{\placeholderfigure}[2]{
  \begin{figure}[h]
    \centering
    \fbox{
      \begin{minipage}[c][5cm][c]{0.82\linewidth}
        \centering
        #1
      \end{minipage}
    }
    \caption{#2}
  \end{figure}
}

\begin{document}

\begin{titlepage}
  \centering
  \vspace*{1.2cm}

  {\Large F.CSM314 -- Програм хангамжийн архитектур\par}
  \vspace{1.5cm}

  {\LARGE\bfseries Багийн ажил №8\par}
  \vspace{0.4cm}
  {\Large\bfseries User \& Access Management модулын хөгжүүлэлтийн тайлан\par}

  \vspace{2cm}

  \begin{tabular}{rl}
    Систем: & Internet Cafe Management System \\
    Модуль: & User \& Access Management \\
    Баг: & \textit{[Багийн нэр]} \\
    Гишүүд: & \textit{[Гишүүдийн нэрс]} \\
    Огноо: & \textit{[Огноо]} \\
  \end{tabular}

  \vfill
  {\large Улаанбаатар хот\par}
\end{titlepage}

\tableofcontents
\newpage

\section{Модулын тодорхойлолт}

Манай баг интернет кафены удирдлагын системээс \textbf{User \& Access Management} модулийг сонгож хөгжүүлсэн. Энэ модуль нь системд хэн нэвтрэх, ямар эрхтэй байх, хэрэглэгчийн төлөв ямар байх, мөн аюулгүй байдлын чухал үйлдлүүд хэрхэн бүртгэгдэхийг хариуцна.

Интернет кафе шиг олон хэрэглэгчтэй системд зөвхөн ``login хийх'' хэсэг хангалтгүй. Админ хэрэглэгчийн жагсаалтыг харах, хэрэглэгчийг идэвхгүй болгох, эрх өөрчлөх, нууц үг сэргээх, гарсан үйлдлийг дараа нь шалгах боломж хэрэгтэй. Тиймээс энэ модулийг authentication, authorization, хэрэглэгчийн удирдлага, activity log гэсэн үндсэн хэсгүүдтэйгээр хийсэн.

Компонент диаграм дээрх \textbf{User Service Port} нь энэ модулийн гадна талтай харьцах интерфэйс гэж ойлгогдож байгаа. Код дээр үүнийг тусдаа класс болгож бичээгүй, харин backend-ийн REST API endpoint-уудаар хэрэгжүүлсэн. Өөрөөр хэлбэл \texttt{/api/auth} болон \texttt{/api/users} route-ууд нь тухайн port-ийн бодит хэрэгжилт болж ажиллаж байна.

\subsection{Модулын хамрах хүрээ}

Энэ модуль дараах үйлдлүүдийг хийж чадна.

\begin{itemize}
  \item шинэ хэрэглэгч бүртгэх;
  \item email болон нууц үгээр login хийх;
  \item JWT token ашиглан хамгаалагдсан API руу хандах;
  \item logout хийх үед token-ийг сервер талд хүчингүй болгох;
  \item хэрэглэгч өөрийн profile мэдээллийг харах, засах;
  \item админ хэрэглэгчдийн жагсаалтыг харах, хайх, шүүх;
  \item админ хэрэглэгчийн role болон status өөрчлөх;
  \item нууц үг сэргээх token үүсгэх, token-оор шинэ нууц үг тохируулах;
  \item login, failed login, logout, password reset, role/status change зэрэг үйлдлийг activity log-д хадгалах.
\end{itemize}

\subsection{Үндсэн шаардлагууд}

Багийн ажил №8-ийн зааварт модулийг хөгжүүлэхдээ үндсэн шаардлагуудаа хангаж, аюулгүй байдлын тактикуудыг хэрэгжүүлэхийг шаардсан. Үүний дагуу бид дараах шаардлагуудыг авсан.

\begin{longtable}{p{0.25\linewidth} p{0.65\linewidth}}
\toprule
\textbf{Шаардлага} & \textbf{Тайлбар} \\
\midrule
Хэрэглэгч бүртгэл & Системд шинэ user account үүсгэх боломжтой байх. \\
Нэвтрэлт & Email, password шалгаж, зөв бол JWT token буцаах. \\
Эрхийн хяналт & ADMIN болон USER role ялгаж, admin endpoint-уудыг зөвхөн ADMIN-д зөвшөөрөх. \\
Profile удирдлага & Хэрэглэгч өөрийн нэр, email зэрэг мэдээллээ харах, засах. \\
User management & Админ user жагсаалт харах, role/status өөрчлөх. \\
Password reset & Нууц үг мартсан үед token үүсгэж, шинэ password тохируулах. \\
Activity log & Аюулгүй байдалтай холбоотой гол үйлдлүүдийг бүртгэх. \\
Туршилт & Security flow-уудыг backend test-ээр шалгах. \\
\bottomrule
\end{longtable}

\section{Ашигласан технологи}

Модулийг backend болон frontend гэж хоёр folder-т салгаж хийсэн. Ингэснээр API болон хэрэглэгчийн дэлгэцийн код тусдаа, ойлгомжтой болсон.

\subsection{Backend}

Backend нь Node.js дээр Express.js ашиглаж хийгдсэн. Өгөгдлийг SQLite database-д хадгалсан. Энэ нь сургалтын ажилд хангалттай хөнгөн, setup хийхэд амар, мөн хэрэглэгч, token, log зэрэг хүснэгтүүдийг цэгцтэй хадгалах боломж өгсөн.

\begin{itemize}
  \item \textbf{Node.js, Express.js} -- REST API үүсгэхэд ашигласан.
  \item \textbf{SQLite} -- хэрэглэгч, activity log, revoked token, password reset token хадгалсан.
  \item \textbf{jsonwebtoken} -- JWT access token үүсгэх, шалгах.
  \item \textbf{bcryptjs} -- password hash хийх.
  \item \textbf{express-validator} -- request body, query, parameter шалгах.
  \item \textbf{helmet} -- HTTP security header-ууд нэмэх.
  \item \textbf{express-rate-limit} -- production орчинд login attempt хязгаарлах.
  \item \textbf{dotenv} -- \texttt{.env} тохиргоо унших.
\end{itemize}

\subsection{Frontend}

Frontend нь React болон Vite ашиглаж хийгдсэн. Нэг дэлгэц дээр login/register, profile, users, activity log, password reset хэсгүүдийг харуулдаг. Admin login хийсэн үед admin-д зориулсан Users болон Activity log цэс харагдана.

\begin{itemize}
  \item \textbf{React} -- component болон state удирдах.
  \item \textbf{Vite} -- frontend dev server болон build хийх.
  \item \textbf{HTML/CSS/JavaScript} -- UI бүтэц, style, API request.
\end{itemize}

\section{Системийн үйл ажиллагаа}

Системийн ерөнхий ажиллагаа энгийн REST API загвараар явна. Browser дээрх React frontend backend рүү HTTP request явуулна. Backend тухайн request-ийг validate хийж, шаардлагатай бол JWT token шалгаад SQLite database-тэй ажиллана.

\begin{center}
\begin{tabular}{c}
Frontend / React UI \\
$\downarrow$ HTTP request \\
Express API \\
$\downarrow$ validation, auth, role check \\
SQLite database
\end{tabular}
\end{center}

\subsection{API endpoint-ууд}

Гол endpoint-ууд:

\begin{lstlisting}[style=api]
POST  /api/auth/register
POST  /api/auth/login
GET   /api/auth/me
POST  /api/auth/logout
POST  /api/auth/forgot-password
POST  /api/auth/reset-password

PATCH /api/users/me
GET   /api/users
PATCH /api/users/:id/status
PATCH /api/users/:id/role
GET   /api/users/activity-logs
\end{lstlisting}

\subsection{Database бүтэц}

Backend эхлэх үед шаардлагатай хүснэгтүүдийг автоматаар үүсгэнэ.

\begin{itemize}
  \item \texttt{users} -- хэрэглэгчийн үндсэн мэдээлэл, password hash, role, status, \texttt{token\_version}.
  \item \texttt{activity\_logs} -- login, logout, password reset зэрэг үйлдлийн бүртгэл.
  \item \texttt{revoked\_tokens} -- logout хийсэн JWT-ийн \texttt{jti}-г хадгалж дахин ашиглуулахаас хамгаална.
  \item \texttt{password\_reset\_tokens} -- reset token-ий hash, expire хугацаа, ашигласан эсэхийг хадгална.
\end{itemize}

\subsection{Хэрэглэгчийн энгийн flow}

Энгийн USER role-той хэрэглэгч дараах байдлаар ажиллана.

\begin{enumerate}
  \item Register хэсгээр account үүсгэнэ.
  \item Email, password ашиглан login хийнэ.
  \item Login амжилттай бол frontend token хадгалж, дараагийн protected request-д ашиглана.
  \item Profile хэсгээс өөрийн мэдээллийг харж, засна.
  \item Password reset хэсгээр reset token үүсгэж, password сольж болно.
  \item Logout хийхэд token revoked болж, хуучин token дахин ажиллахгүй.
\end{enumerate}

\placeholderfigure{Login болон Register дэлгэцийн screenshot-оо энд оруулна.}{Login/Register дэлгэц}

\subsection{Админы flow}

ADMIN role-той хэрэглэгч системийн хэрэглэгчдийг удирдана. Admin login хийсний дараа Users болон Activity log цэс гарч ирнэ.

\begin{enumerate}
  \item Users хэсэгт хэрэглэгчдийн жагсаалт гарна.
  \item Search, role, status filter ашиглаж хэрэглэгч хайна.
  \item Хэрэглэгчийн status-ийг ACTIVE эсвэл INACTIVE болгож өөрчилнө.
  \item Хэрэглэгчийн role-ийг USER эсвэл ADMIN болгож өөрчилнө.
  \item Activity log хэсгээс системд болсон чухал үйлдлүүдийг шалгана.
\end{enumerate}

\placeholderfigure{Admin Users дэлгэцийн screenshot-оо энд оруулна.}{Admin хэрэглэгчийн удирдлагын дэлгэц}

\placeholderfigure{Activity log дэлгэцийн screenshot-оо энд оруулна.}{Activity log дэлгэц}

\section{Аюулгүй байдлын тактикуудын хэрэгжилт}

Энэ модуль дээр аюулгүй байдлыг зөвхөн нэг функцээр шийдэх биш, хэд хэдэн жижиг хамгаалалтыг хамтад нь ашигласан. Доорх тактикууд хамгийн чухал хэсгүүд нь байсан.

\subsection{Authentication}

Login хийх үед хэрэглэгчийн email болон password шалгагдана. Password database-д plain text хэлбэрээр хадгалагдахгүй, зөвхөн bcrypt hash хадгалагдана. Login амжилттай бол backend JWT token үүсгэж frontend рүү буцаана.

JWT дотор user id, role, token version, мөн \texttt{jti} гэсэн unique id ордог. Энэ нь дараа нь token revoke болон session invalidation хийхэд хэрэг болсон.

\subsection{Authorization}

Зарим endpoint зөвхөн ADMIN role-той хэрэглэгчид зориулагдсан. Жишээ нь user list харах, role/status өөрчлөх, activity log харах зэрэг үйлдлүүдэд эхлээд JWT шалгаад, дараа нь role шалгадаг.

Ингэснээр энгийн USER өөрийн profile-оо харж засаж болох ч бусад хэрэглэгчийн эрх, төлөвийг өөрчлөх боломжгүй.

\subsection{Password protection}

Нууц үгийг bcrypt ашиглан hash хийж хадгалсан. Энэ нь database алдагдсан ч хэрэглэгчийн password шууд харагдахгүй гэсэн үг. Мөн password reset хийх үед хуучин password-оор login хийх боломжгүй болж, шинэ password л хүчинтэй үлдэнэ.

\subsection{Token revoke ба session invalidation}

JWT token нь ердийн үед stateless байдаг. Гэхдээ logout хийсний дараа хуучин token ажилласаар байвал бодит системд асуудал болно. Тиймээс logout хийх үед token-ийн \texttt{jti}-г \texttt{revoked\_tokens} хүснэгтэд хадгалсан. Дараа нь protected endpoint рүү хандах бүрд token revoked эсэхийг шалгана.

Мөн user-ийн role, status, password өөрчлөгдөхөд \texttt{token\_version} нэмэгдэнэ. JWT доторх token version database дээрх утгатай таарахгүй бол token хүчингүй болно. Энэ нь хэрэглэгчийн эрх өөрчлөгдсөн ч хуучин token-оор хуучин эрхээ ашиглахаас хамгаална.

\subsection{Password reset}

Password reset token санамсаргүй байдлаар үүснэ. Database-д token өөрөө биш, hash нь хадгалагдана. Token нь хугацаатай, нэг удаа ашиглагддаг. Ашиглагдсан эсвэл хугацаа дууссан token-оор password reset хийх боломжгүй.

Demo/development үед reset token response-д харагдаж байгаа. Production орчинд үүнийг email эсвэл notification service-ээр явуулах нь зөв.

\subsection{Input validation}

Request бүр дээр email, password, full name, role, status, page, limit, user id зэрэг утгуудыг шалгадаг. Ингэснээр буруу format-тай өгөгдөл database рүү орохоос өмнө буцаагдана.

\subsection{Rate limiting}

Production орчинд login endpoint дээр rate limit ажиллана. Энэ нь богино хугацаанд олон password оролдлого хийхээс хамгаална. Development үед frontend туршихад саад болохгүй байхаар rate limit-ийг production орчинд л ажиллуулахаар тохируулсан.

\subsection{Audit log}

Security талаасаа дараа нь шалгаж болох бүртгэл хэрэгтэй. Тиймээс дараах үйлдлүүдийг activity log-д хадгалж байгаа.

\begin{itemize}
  \item REGISTER
  \item LOGIN
  \item FAILED\_LOGIN
  \item LOGOUT
  \item PASSWORD\_RESET\_REQUESTED
  \item PASSWORD\_RESET
  \item STATUS\_CHANGE
  \item ROLE\_CHANGE
\end{itemize}

Activity log нь админд систем дээр юу болсон, хэзээ болсон, ямар user-тэй холбоотой байсныг харах боломж өгнө.

\subsection{Admin хамгаалалт}

Админ эрхийн өөрчлөлт дээр нэг нэмэлт хамгаалалт хийсэн. Системд дор хаяж нэг active admin үлдэх ёстой. Сүүлийн active admin-ийг inactive болгох эсвэл admin role-оос буулгахыг backend зөвшөөрөхгүй. Энэ хамгаалалт байхгүй бол системд admin эрхтэй хүн үлдэхгүй болох эрсдэлтэй.

\section{Туршилт}

Backend дээр Node.js-ийн test runner ашиглаж integration test бичсэн. Test нь зөвхөн нэг endpoint шалгаад зогсохгүй, module-ийн гол security flow-уудыг дарааллаар нь шалгадаг.

Туршсан гол зүйлс:

\begin{itemize}
  \item хэрэглэгч бүртгэх;
  \item admin login хийх;
  \item user жагсаалт pagination болон role filter-тэй ажиллах;
  \item logout хийсэн token дахин ашиглагдахгүй байх;
  \item inactive user зөв password-той байсан ч login хийж чадахгүй байх;
  \item status өөрчлөгдөхөд хуучин token хүчингүй болох;
  \item password reset хийсний дараа хуучин password ажиллахгүй болох;
  \item user-ийг admin болгосны дараа admin endpoint ашиглаж болох;
  \item activity log уншигдах.
\end{itemize}

Ажиллуулсан command:

\begin{lstlisting}[style=api]
npm run backend:test
npm run frontend:build
\end{lstlisting}

Сүүлд шалгахад backend test болон frontend build хоёулаа амжилттай болсон.

\section{GitHub дээр байршуулах}

Эх кодыг GitHub repository дээр байршуулна. Repository-д backend, frontend хоёр тусдаа folder-той байна.

\begin{lstlisting}[style=api]
backend/   Express + SQLite API
frontend/  React + Vite frontend
\end{lstlisting}

GitHub холбоос:

\begin{center}
\texttt{https://github.com/username/repository}
\end{center}

\section{Дүгнэлт}

Энэ ажлаар бид интернет кафены системийн хэрэглэгч болон эрхийн удирдлагын модулийг ажилладаг байдлаар хөгжүүлсэн. Эхэндээ зөвхөн login/register хэсэг шиг харагдаж байсан ч бодитоор хийхэд token revoke, role check, password reset, activity log, inactive account, last admin protection зэрэг олон жижиг шийдэл хэрэгтэй болсон.

Компонент диаграм дээр байсан User \& Access Management хэсгийг backend API болон frontend дэлгэцээр бодит болгосон. Диаграмын User Service Port нь код дээр REST endpoint-уудаар илэрхийлэгдэж байгаа. Модуль нь хэрэглэгчийн үндсэн үйлдлүүдийг хийж чадна, админ удирдлагатай, мөн аюулгүй байдлын хэд хэдэн тактик хэрэгжүүлсэн.

Цаашид production түвшинд ашиглах бол password reset token-ийг email service-ээр явуулах, HTTPS тохируулах, frontend дээр илүү сайн form validation нэмэх, мөн refresh token эсвэл session management-ийг илүү боловсронгуй болгох боломжтой. Гэхдээ багийн ажил №8-ийн хүрээнд сонгосон модулийн үндсэн шаардлага болон security tactic-уудыг хангасан гэж үзэж байна.

\end{document}
